\documentclass[11pt]{article}
\usepackage{amsmath,amssymb}
\usepackage{physics}
\usepackage[margin=1in]{geometry}
\usepackage{enumitem}

\title{ATPL 2025 --- Exercises: The Output Problem and Amplitude
Amplification}
\author{Lecture 1: Why do we need hybrid programming models?}
\date{}

\begin{document}
\maketitle

Background: Nielsen \& Chuang, Chapter 6 (Grover's algorithm); Brassard,
H{\o}yer, Mosca \& Tapp, \emph{Quantum Amplitude Amplification and
Estimation} (2002), for the general (non-Grover-specific) treatment used
below.

\begin{enumerate}[label=\textbf{Exercise \arabic*.}, leftmargin=2.2cm]

\item \textbf{The cost of reading out a full amplitude vector.}
Suppose an $n$-qubit state $\ket\Psi = \sum_i a_i \ket i$ has a roughly
flat distribution, $|a_i|^2 \approx 2^{-n}$ for all $i$. Using a basic
concentration-of-measure argument (e.g.\ the standard error of a
binomial proportion estimate from $S$ samples is $\approx
\sqrt{p(1-p)/S}$), estimate how many repeated runs $S$ are needed to
estimate \emph{every} $|a_i|^2$ to additive precision $\varepsilon$ with
reasonable confidence. Confirm this matches the $\mathcal O(2^n)$
(with an $\varepsilon$-dependent constant) claim on the lecture's Output
Problem slide.

\item \textbf{Amplitude amplification as a 2D rotation.}
Let $A$ be a state-preparation circuit with
$A\ket0 = \sin\theta\ket{\mathrm{good}} + \cos\theta\ket{\mathrm{bad}}$,
where $\ket{\mathrm{good}},\ket{\mathrm{bad}}$ span a 2-dimensional
subspace and $\theta = \arcsin\sqrt p$ for success probability $p$.
\begin{enumerate}[label=(\alph*)]
\item Define the oracle reflection
$S_{\mathrm{good}} = I - 2\ket{\mathrm{good}}\!\bra{\mathrm{good}}$
and the amplitude reflection
$S_0 = 2A\ket0\!\bra0A^\dagger - I$.
Show that the amplitude-amplification iterate
$G = -S_0 S_{\mathrm{good}}$
acts as a rotation by angle $2\theta$ within the
$\{\ket{\mathrm{good}},\ket{\mathrm{bad}}\}$ plane.
\item Show that after $k$ applications of $G$ to $A\ket0$, the success
probability is $\sin^2\!\big((2k{+}1)\theta\big)$, and derive the
(approximately) optimal number of iterations
$k^\ast \approx \tfrac\pi4\sqrt{1/p} - \tfrac12$
that maximizes this probability.
\end{enumerate}

\item \textbf{Small worked example.}
Take $N=4$ (i.e.\ $n=2$ qubits), with exactly one marked item, so
$p = 1/4$ and $\theta = \pi/6$.
\begin{enumerate}[label=(\alph*)]
\item Compute $\theta$ exactly, and find the optimal integer $k^\ast$
from the formula in Exercise 2b.
\item By explicit $2\times2$ rotation-matrix multiplication (in the
$\{\mathrm{good},\mathrm{bad}\}$ basis), compute the success
probability after $k=0,1,2$ iterations, and confirm that it is maximized
(in fact reaches exactly $1$) at $k=k^\ast$.
\end{enumerate}

\item \textbf{Quadratic speedup for the Output Problem.}
Suppose you don't know in advance which basis state is "the answer,"
but you can efficiently check correctness classically (an oracle).
Using amplitude amplification, explain why you can find a correct
answer using $\mathcal O\!\left(\sqrt{2^n/M}\right)$ applications of a
state-preparation circuit $A$ (where $M$ is the number of correct
answers among $2^n$ possibilities), rather than the naive
$\mathcal O(2^n/M)$ implied by repeated classical-style sampling from
the Output Problem slide.
\begin{enumerate}[label=(\alph*)]
\item Is this an exponential or a polynomial (quadratic) improvement in
$n$?
\item How does this relate to why Grover's algorithm alone is not
believed to be enough to place NP-complete problems inside BQP (recall
the complexity-class picture from the lecture)?
\end{enumerate}

\end{enumerate}

\end{document}
